\subsubsection*{مراحل داکرایز کردن}

\subsubsection*{مرحله اول: \lr{Builder Stage}}

\begin{itemize}
    \item \textbf{\lr{FROM python:3.10-slim AS builder}:}
    ایجاد یک مرحلهٔ مستقل با نام \lr{builder} برای ایزوله کردن فرایند ساخت و کامپایل وابستگی‌ها.

    \item \textbf{\lr{ENV PYTHONDONTWRITEBYTECODE=1}:}
    جلوگیری از تولید فایل‌های \lr{.pyc}.

    \item \textbf{\lr{ENV PYTHONUNBUFFERED=1}:}
    ارسال مستقیم خروجی‌های \lr{stdout} و \lr{stderr} به لاگ‌های داکر بدون \lr{buffer} شدن.

    \item \textbf{\lr{ARG MIRROR\_URL=""}:}
    تعریف یک متغیر مخصوص زمان \lr{build} برای تغییر منبع دانلود بسته‌ها.

    \item \textbf{\lr{RUN --mount=type=cache,target=/root/.cache/pip}:}
    استفاده از قابلیت \lr{BuildKit Cache Mount} برای کش کردن بسته‌های دانلودشده توسط \lr{pip} و کاهش زمان بیلدهای بعدی.

    \item \textbf{\lr{--prefix=/install}:}
    نصب وابستگی‌ها در مسیر ایزولهٔ \lr{/install} تا بتوان آن‌ها را به مرحلهٔ نهایی منتقل کرد.
\end{itemize}

\subsubsection*{مرحله دوم: \lr{Runner Stage (Final Image)}}

\begin{itemize}
    \item \textbf{\lr{FROM python:3.10-slim AS runner}:}
    شروع یک مرحلهٔ جدید و سبک بدون ابزارها و کش‌های مرحلهٔ ساخت.

    \item \textbf{\lr{PATH} و \lr{PYTHONPATH}:}
    افزودن مسیر بسته‌های نصب‌شده به متغیرهای محیطی.

    \item \textbf{\lr{RUN addgroup ... \&\& adduser ...}:}
    ایجاد یک کاربر \lr{Non\mbox{-}root} برای افزایش امنیت و جلوگیری از حملات \lr{Container Breakout}.

    \item \textbf{\lr{COPY --from=builder /install /install}:}
    انتقال وابستگی‌های نصب‌شده بدون انتقال ابزارهای مرحلهٔ ساخت.

    \item \textbf{\lr{RUN chown -R appuser:appgroup /app}:}
    واگذاری مالکیت فایل‌های پروژه به کاربر غیرریشه.

    \item \textbf{\lr{USER appuser}:}
    اجرای برنامه با کاربر غیرریشه.
\end{itemize}

\subsubsection*{نحوهٔ Build و اجرا}

ساخت ایمیج با فعال بودن \lr{BuildKit}:

\begin{latin}
\begin{lstlisting}[language=bash]
DOCKER_BUILDKIT=1 docker build \
  --build-arg MIRROR_URL=https://pypi.tuna.tsinghua.edu.cn/simple \
  -t my-python-app:v1 .
\end{lstlisting}
\end{latin}

اجرای کانتینر:

\begin{latin}
\begin{lstlisting}[language=bash]
docker run --rm my-python-app:v1
\end{lstlisting}
\end{latin}

\subsubsection*{نقش کلیدواژه‌ها}

\begin{description}
    \item[\lr{ENV}]
    متغیر محیطی تعریف می‌کند که هم در زمان \lr{build} و هم در زمان اجرای کانتینر در دسترس است و در \lr{metadata} ایمیج ذخیره می‌شود. بنابراین محل مناسبی برای نگهداری اطلاعات محرمانه نیست.

    \item[\lr{WORKDIR}]
    پوشهٔ کاری دستورات بعدی مانند \lr{RUN}، \lr{COPY} و \lr{CMD} را تعیین می‌کند و در صورت نیاز آن را ایجاد می‌کند.

    \item[\lr{AS}]
    به یک مرحلهٔ ساخت نام می‌دهد تا بتوان از آن در دستورهایی مانند \lr{COPY --from} یا \lr{--target} استفاده کرد.
\end{description}

\subsubsection*{\lr{ADD} در برابر \lr{COPY}}

هر دو فایل‌ها را از \lr{build context} به داخل ایمیج منتقل می‌کنند، اما \lr{ADD} دو قابلیت اضافه دارد:

\begin{itemize}
    \item استخراج خودکار آرشیوهای محلی؛
    \item دانلود فایل از یک \lr{URL} (و در نسخه‌های جدید \lr{BuildKit}، کلون کردن مخزن گیت).
\end{itemize}

به همین دلیل معمولاً استفاده از \lr{COPY} توصیه می‌شود و فقط در صورت نیاز به استخراج خودکار آرشیو از \lr{ADD} استفاده می‌شود.

\subsubsection*{\lr{EXPOSE} در برابر \lr{Port\mbox{-}Mapping}}

\lr{EXPOSE 8000} فقط اعلام می‌کند که برنامه روی این پورت گوش می‌دهد و هیچ پورتی را منتشر نمی‌کند.

در مقابل، \lr{docker run -p 8000:8000} یا کلید \lr{ports} در \lr{Docker Compose} باعث ایجاد نگاشت واقعی پورت روی میزبان می‌شود.

\subsubsection*{\lr{ARG} در برابر \lr{ENV}}

\begin{center}
\renewcommand{\arraystretch}{1.3}
\begin{tabular}{|p{4cm}|p{4.5cm}|p{5cm}|}
\hline
 & \centering\lr{ARG} & \centering\arraybackslash\lr{ENV} \\
\hline
دامنهٔ حیات & فقط زمان \lr{build} & زمان \lr{build} و زمان اجرا \\
\hline
تعیین مقدار & \lr{--build-arg} & \lr{-e}، \lr{env\_file} یا مقدار پیش‌فرض \\
\hline
حضور در ایمیج نهایی & خیر (اما در \lr{docker history}) & بله، در \lr{metadata} \\
\hline
دیده شدن در برنامه & خیر & بله \\
\hline
دامنهٔ اثر & فقط همان مرحله & به مراحل بعدی ارث می‌رسد \\
\hline
\end{tabular}
\end{center}

البته که هیچ‌کدام برای نگهداری \lr{secret} مناسب نیستند.

\subsubsection*{\lr{CMD} در برابر \lr{ENTRYPOINT}}

\lr{ENTRYPOINT} مشخص می‌کند کانتینر چه فرایندی را اجرا می‌کند و معمولاً ثابت در نظر گرفته می‌شود؛ \lr{CMD} آرگومان‌های پیش‌فرضی است که به آن فرایند داده می‌شود و کاربر می‌تواند هنگام \lr{docker run} آن را بازنویسی کند. وقتی هر دو با هم بیایند، دستور نهایی از الحاق این دو ساخته می‌شود:

\begin{latin}
\begin{lstlisting}[language=Docker]
ENTRYPOINT ["uvicorn"]
CMD ["app.main:app", "--host", "0.0.0.0", "--port", "8000"]
\end{lstlisting}
\end{latin}

\begin{itemize}
    \item \lr{docker run img}
    $\leftarrow$
    \lr{uvicorn app.main:app --host 0.0.0.0 --port 8000}
    (چیزی داده نشده، پس \lr{CMD} پیش‌فرض به \lr{ENTRYPOINT} الحاق می‌شود.)

    \item \lr{docker run img --version}
    $\leftarrow$
    \lr{uvicorn --version}
    (آرگومان داده‌شده فقط جای \lr{CMD} را می‌گیرد، \lr{ENTRYPOINT} ثابت می‌ماند.)

    \item تغییر خود \lr{ENTRYPOINT} فقط با فلگ
    \lr{docker run --entrypoint ...}
    امکان‌پذیر است.
\end{itemize}

اگر فقط \lr{CMD} تعریف شود (بدون \lr{ENTRYPOINT})، هر آرگومانی که در \lr{docker run} داده شود کل دستور را جایگزین می‌کند؛ این رفتار برای ایمیج‌های چندمنظوره مناسب است، ولی برای یک سرویس مشخص، ترکیب \lr{ENTRYPOINT} به‌همراه \lr{CMD} رفتاری ثابت‌تر و قابل‌پیش‌بینی‌تر می‌دهد.

نکتهٔ مهم دیگر این‌که هر دو باید به‌صورت \lr{exec} (آرایهٔ \lr{JSON}، مانند مثال بالا) نوشته شوند، نه به‌صورت \lr{shell}. در قالب \lr{shell} (مثلاً \lr{CMD python main.py})، فرایند اصلی برنامه زیرِ \lr{/bin/sh} اجرا می‌شود و خودش \lr{PID 1} نیست؛ در نتیجه سیگنال \lr{SIGTERM} که \lr{docker stop} می‌فرستد به آن نمی‌رسد و کانتینر به‌جای خاموشی تمیز، تا رسیدن \lr{timeout} معطل می‌ماند و سپس با \lr{SIGKILL} بسته می‌شود.

\subsubsection*{ساخت و اجرای کانتینر از روی \lr{Dockerfile}}

\begin{latin}
\begin{lstlisting}[language=bash]
# build
docker build -t my-python-app:v1 .

# run
docker run --rm my-python-app:v1

# verify
docker run --rm my-python-app:v1 whoami

# check size and layers
docker image inspect my-python-app:v1 --format '{{.Size}}'
docker history my-python-app:v1
\end{lstlisting}
\end{latin}

از آن‌جا که این برنامه یک اسکریپت یک‌بارمصرف است، نه یک سرویس شنونده روی پورت، نیازی به \lr{-p}، \lr{EXPOSE} یا \lr{docker compose} ندارد؛ کانتینر پس از اجرای \lr{main.py} و چاپ خروجی به‌صورت خودکار خاتمه می‌یابد و به‌لطف \lr{--rm} حذف هم می‌شود.